% ============================================================================
%  The Bitcoin Block Space Problem:
%  Does the Fee Market Internalize Long-Term Storage Costs?
%  Working Paper v2.1.0 (model-spec.json v2.0.1) — LaTeX source
%
%  CONVERSION STATUS (honest):
%   - Abstract: FULLY CONVERTED from working-paper.md (verbatim).
%   - Sections 1-10: structure + all tables + key prose converted; some
%     long prose passages are condensed but faithful (see comments).
%   - References: converted from the paper's numbered list.
%   - NOT YET CONVERTED / STUBS: figures (none in the .md), the
%     verification-appendix content (kept as companion doc reference),
%     and full verbatim prose of every section (the .md remains the
%     canonical text). This source compiles; content parity with the .md
%     should be checked by diff before submission.
%
%  2026-08-03: subsection 7.1 (falsifiability) added to this source.
%  2026-08-03: NOTE working-paper.md sections 11-12 (deep questions +
%   cross-chain) are NOT YET CONVERTED to LaTeX (open item, see repo handoff).
%
%  TOOLCHAIN: pdflatex (not present on the dev machine at conversion
%  time — compile on any TeX installation:  pdflatex working-paper.tex)
%
%  AUTHOR: Prateek Poswal (identity choice pending — see
%  research/author-identity.md; swap the \author block once ratified).
% ============================================================================
\documentclass[11pt]{article}
\usepackage[margin=1in]{geometry}
\usepackage{amsmath,amssymb}
\usepackage{booktabs}
\usepackage{array}
\usepackage[T1]{fontenc}
\usepackage[utf8]{inputenc}
\usepackage[hidelinks]{hyperref}
\usepackage{microtype}
\usepackage{enumitem}

\title{\textbf{Storage Cost Internalization in Bitcoin's Fee Market}\\[0.4em] \large The Bitcoin Block Space Problem (program subtitle)}
\author{Prateek Poswal\\ \small Independent Researcher (Bitcoin Sahi Research)}
\date{Working Paper v2.1.0 \textbullet{} model-spec v2.0.1 \textbullet{} 2026-08-02}

\begin{document}
\maketitle

\begin{abstract}
Bitcoin's fee market allocates scarce block space among competing transactions. This paper asks a distinct empirical question: does the fee market \emph{also} internalize the long-term resource costs imposed by permanently recorded blockchain data? We define a \textbf{Storage Cost Coverage Ratio (SCCR)} --- the ratio of transaction fees paid (USD) to the estimated lifetime storage cost borne by full nodes (USD) --- and measure it against live fee-history data using a \textbf{primary-source node census ($\geq$32{,}000 known addresses from a live Bitcoin Core node)}. Across 156 sampled blocks, the average ratio is $\sim$0.29 (dimensionless): fees cover roughly 29\% of the estimated 10-year replicated storage cost of an average block, and $\sim$99\% of sampled blocks fall below the $1\times$ threshold (a small number of high-fee blocks exceed it). A live re-measure on 2026-08-02 gives $\sim$0.225 (dimensionless) (168 blocks, model-spec v2.0.1); a re-read of the same capture at time of writing (167 blocks) gave 0.2228 with 100\% of blocks below $1\times$ --- the ratio moves with the fee market. We reconcile two cost models that previously disagreed by $16.4\times$ (dimensionless), document the correction transparently (model-spec.json v2.0.0), and bound the result against parameter uncertainty (joint Monte Carlo: 99.8\% of draws below $1\times$ under the old N band; $\sim$99\% at the real census). We present the framework as reproducible, falsifiable research --- not a claim that Bitcoin is broken.
\end{abstract}

\noindent\textbf{Hypothesis:} Bitcoin's fee market efficiently allocates scarce block space, but may not fully internalize every long-lived resource cost created by confirmed transactions.

\section{Introduction: Congestion Pricing vs.\ Permanence Cost}\label{sec:1}
The one-sentence frame (full treatment in \texttt{problem\_statement.md}):
\begin{quote}
\textbf{Bitcoin's fee market prices competition for inclusion in the next block, but it does not explicitly price the long-term resource costs of permanently recorded blockchain data.}
\end{quote}
The fee market solves a short-term optimization problem extremely well: during congestion, higher fees win block space. This is \textbf{congestion pricing}, with a horizon of roughly one block ($\sim$10~min). Storage cost is structurally different:
\begin{itemize}[nosep]
  \item \textbf{One-time payment} (the fee, USD/block) vs.\ \textbf{recurring cost} (long-term replicated storage across the network, USD/(node$\cdot$yr))
  \item \textbf{Payer chooses} to pay vs.\ \textbf{node operators bear} the cost involuntarily
  \item $\sim$10-min horizon (next block) vs.\ \textbf{indefinite horizon} (data persists in blockchain history)
\end{itemize}
\textbf{Scope:} this paper measures the \textbf{storage leg} of a broader resource-pricing question. Bandwidth, validation, and UTXO-maintenance legs are future work (see \S\ref{sec:7}).

\section{What We're NOT Saying}\label{sec:2}
\begin{itemize}[nosep]
  \item $\times$ Not proposing a fix
  \item $\times$ Not claiming the externality is economically significant at current volumes (it may not be)
  \item $\times$ Not claiming the SegWit discount was a mistake (it solved transaction malleability)
  \item $\times$ \textbf{Not arguing that Bitcoin is `broken'} --- the fee market solves block-space allocation well; whether it internalizes long-lived resource costs is an empirical question
  \item $\checkmark$ Just framing the question clearly and measurably
\end{itemize}

\section{The SegWit Pricing Structure}\label{sec:3}
SegWit (BIP~141, activated August 2017) introduced block weight: witness data counts \textbf{1 weight unit (WU) per byte}, non-witness data \textbf{4~WU per byte}. This \textbf{4:1 (dimensionless) discount} reduces the marginal block-space cost of witness-resident data relative to non-witness data. Years later, inscription protocols (Ordinals, BRC-20, Runes) took advantage of that pricing structure. Two precise caveats (full treatment in \texttt{bip141\_analysis.md}):
\begin{enumerate}[nosep]
  \item The discount applies to \textbf{all} witness data --- SegWit financial transactions benefit identically. The ``accident'' is that it also subsidizes data-bearing constructions, not that SegWit was designed for them.
  \item Whether the discount appropriately reflects the long-term resource costs of the data it makes cheaper is \textbf{part of the research question}, not a foregone conclusion.
\end{enumerate}

\section{The Model}\label{sec:4}
\subsection{Canonical specification}\label{sec:4.1}
All quantities, units, and formulas live in \texttt{research/model-spec.json} (v2.0.1) --- a single canonical source that every script imports. No script redefines a model constant.
\begin{center}
\renewcommand{\arraystretch}{1.25}
\begin{tabular}{lllll}
\toprule
Quantity & Symbol & Units & Value & Source \\
\midrule
Annual node cost & $C$ & USD/yr & 925 & \texttt{utxo\_cost\_model} component sum (924.35) rounded \\
Replication factor & $N$ & nodes & 32{,}000 & primary-source census (agent-25, \texttt{getnodeaddresses}, $\geq$32K) \\
Storage horizon & $T$ & yr & 10 & assumption (archival retention) \\
Average block size & $B_{\mathrm{block}}$ & bytes & 1{,}500{,}000 & \texttt{block\_stats} / \texttt{fee\_history} \\
Blocks per year & $R_{\mathrm{blocks}}$ & blocks/yr & 52{,}596 & $365.25 \times 24 \times 6$ \\
Total block bytes/yr & $B_{\mathrm{all,yr}}$ & bytes/yr & $7.8894\times10^{10}$ & $B_{\mathrm{block}} \times R_{\mathrm{blocks}}$ \\
\textbf{Cost per byte per year} & \textbf{$c_b$} & USD/(byte$\cdot$yr) & \textbf{$1.17246\times10^{-8}$} & $C / B_{\mathrm{all,yr}}$ (block-average) \\
Lifetime storage cost/node/block & $L$ & USD/block & 0.175869 & $c_b \times B_{\mathrm{block}} \times T$ \\
Network lifetime cost/block & $L_{\mathrm{net}}$ & USD/block & 5{,}628 & $L \times N$ ($=32$K) \\
\textbf{Storage Cost Coverage Ratio} & \textbf{SCCR} & dimensionless & $\sim$0.29 (156 blocks, N=32K, 2026-08-01); 0.2228 (167 blocks, 2026-08-02) & $\mathrm{fee}_{\mathrm{USD}} / L_{\mathrm{net}}$ \\
\bottomrule
\end{tabular}
\end{center}
\textbf{Design principle:} $c_b$ is \textbf{horizon-free} ($C /$ annual bytes). The horizon $T$ enters \textbf{only} through $L = c_b \times B \times T$. This corrects the v1.0.0 implementation, which applied $T$ twice (see \S\ref{sec:6}).

\subsection{The marginal-inscription attribution (secondary)}\label{sec:4.2}
The inscription-externality branch uses a \textbf{marginal attribution}: $c_{b,\mathrm{insc}} = C / (\text{inscription bytes/yr}) = 1.92573\times10^{-6}$~USD/(byte$\cdot$yr). This differs from $c_b$ by a factor of \textbf{164.4$\times$ (dimensionless)} --- not an error, but two different denominators (all block bytes vs.\ inscription-only bytes). The paper headline uses the block-average $c_b$; the marginal figure appears only in the inscription-externality analysis. The $16.4\times$ (dimensionless) figure reported in earlier versions was this $164\times$ (dimensionless) denominator gap \textbf{divided by} the $10\times$ (dimensionless) bug --- see \S\ref{sec:6}.

\section{Findings}\label{sec:5}
\subsection{The Storage Cost Coverage Ratio}\label{sec:5.1}
Measured from live \texttt{fee\_history} captures, node count $N=32{,}000$ (real census). Two snapshots, dated explicitly:
\begin{center}
\renewcommand{\arraystretch}{1.25}
\begin{tabular}{lcc}
\toprule
Metric & Dated capture (2026-08-01) & Live capture (2026-08-02) \\
\midrule
Blocks sampled & 156 & 167 \\
Average SCCR (dimensionless) & $\sim$0.293 & 0.2228 \\
Min / Max (dimensionless) & $\sim$0.058 / $\sim$1.537 & $\sim$0.058 / $\sim$0.832 \\
Blocks below $1\times$ & 98.7\% (154/156) & 100.0\% \\
\bottomrule
\end{tabular}
\end{center}
The ratio moves with the fee market; the single source of truth is \texttt{research/model-spec.json} (v2.0.1, canonical); all surfaces must read the live value from \texttt{node tools/research/storage-ratio.js}, never a hardcoded figure. The v2.0.0 N=60K-era values (0.1719 / 0.1535) are superseded.

\textbf{Fee-regime + node-count dependence:} the ratio tracks both the fee market and the replication factor. Under the earlier N=60K assumption the average was 0.156--0.172 (dimensionless) with 100\% below $1\times$; at the real census N=32K it rises to $\sim$0.22--0.29 with $\sim$99--100\% below $1\times$ (a few high-fee blocks exceed coverage in the dated capture). The headline is a \emph{distribution over time and parameters}, not a point.

\subsection{The inscription externality (marginal branch)}\label{sec:5.2}
At 100{,}000 inscriptions/month ($\sim$400 bytes each, in witness):
\begin{center}
\begin{tabular}{lc}
\toprule
Metric & Value \\
\midrule
Cost per byte per year (marginal) & $1.92573\times10^{-6}$ USD/(byte$\cdot$yr) \\
Lifetime storage cost per inscription (10 yr) & \$0.00770 USD/inscription \\
New 10-yr storage liabilities created per year at 100K/mo & $\sim$\$9{,}240 USD/yr spread across all nodes \\
Steady-state annual externality (amortized) & $\sim$\$924 USD/yr spread across all nodes \\
\bottomrule
\end{tabular}
\end{center}
\textbf{Unit note (corrected):} \$9{,}240 is the \emph{undiscounted 10-year liability of one year's inscriptions} (USD/yr), not an annual cost. Amortized over the 10-yr horizon it is $\sim$\$924/yr. This was previously labeled ``annual unpriced externality'' --- a T/1 conflation of the same shape as the fixed $10\times$ bug, in label only.

\textbf{Pruning-consistency note:} this externality is structurally real but \textbf{not economically significant at current volumes} --- per-node it is well under \$1/yr. The contribution of this paper is the reproducible ratio framework, not the magnitude.

\subsection{Sensitivity}\label{sec:5.3}
Recomputed at the live capture (2026-08-02, 167 blocks; baseline SCCR $=$ 0.2228 (dimensionless) at N=32K, C=\$925/yr, T=10~yr). The table's parameter ranges bracket the canonical values:
\begin{center}
\renewcommand{\arraystretch}{1.2}
\begin{tabular}{lcl}
\toprule
Parameter & Value & Avg SCCR (dimensionless) \\
\midrule
Node cost, $C$ (USD/yr) & 600 / 925 / 1400 & 0.343 / 0.223 / 0.147 \\
Node count, $N$ (nodes) & 16{,}000 / 32{,}000 / 100{,}000 & 0.446 / 0.223 / 0.071 \\
Storage horizon, $T$ (yr) & 5 / 10 / 15 & 0.446 / 0.223 / 0.149 \\
Avg block size, $B$ (MB) & 1.0 / 1.5 / 2.0 & 0.223 (invariant) \\
\bottomrule
\end{tabular}
\end{center}
\textbf{BTC price (the dominant omitted driver --- verified in independent implementation):} the ratio is linear in price. At the live N=32K baseline (capture price $\approx$ \$63{,}018 USD/BTC): \$30K$\rightarrow$0.106, \$45K$\rightarrow$0.159, \$62.9K$\rightarrow$0.222, \$100K$\rightarrow$0.354, \$200K$\rightarrow$0.707, \$500K$\rightarrow$1.768.

\textbf{Structural property:} SCCR is \textbf{invariant to block size} --- $c_b \propto 1/B$ while $L \propto B$, so $B$ cancels. Caveat: this invariance follows from attributing \emph{all} node cost per byte then multiplying back by bytes --- it reflects that the model is size-agnostic (contains no size-dependent economics), not that size is economically irrelevant.

\subsection{The knife-edge: node count and the strong claim}\label{sec:5.4}
The SCCR is homogeneous of degree 1 in its scale parameters: $\mathrm{SCCR} \propto (\mathrm{fee} \times \mathrm{price}) / (C \times T \times N)$. Two thresholds matter (derived and verified in the independent C implementation):
\begin{itemize}[nosep]
  \item The average SCCR crosses 1.0 at $N \approx$ 10{,}300 nodes (or BTC $\approx$ \$366{,}000) --- dated v2.0.0-era reference values (baseline 0.1719 at N=60K)
  \item The ``100\% of sampled blocks below $1\times$'' claim breaks at $N \approx$ 49{,}200 nodes (or BTC $\approx$ \$76{,}700) --- the highest-fee sampled block (13.75M sats, height 960469) sits at 0.82$\times$ coverage under the old N=60K
\end{itemize}
\textbf{Live recompute (2026-08-02 baseline, SCCR $=$ 0.2228 at N=32K):} the average now inverts at $N \approx$ 7{,}130 nodes or BTC $\approx$ \$283{,}000; the ``100\% below $1\times$'' break on the dated capture is unchanged at $N \approx$ 49{,}200. Both thresholds are reported; model-spec v2.0.1 retains the v2.0.0-era values (10.3K / \$366K).

\textbf{The real census (primary source).} We replaced the 60K assumption with data from a live Bitcoin Core node (\texttt{getnodeaddresses}, agent-25/node-census): \textbf{32{,}000 known addresses} --- the RPC maximum, meaning the node's address manager is saturated at the cap and the true reachable set is \emph{at least} 32K. Independent estimates span $\sim$10K--100K.

\textbf{Consequence of the real N=32{,}000 (recomputed, dated 156-block capture):}
\begin{center}
\begin{tabular}{lcc}
\toprule
Metric & At N=60K (old assumption) & At N=32K (real census) \\
\midrule
Average SCCR (dimensionless) & 0.156--0.172 & \textbf{$\sim$0.293} \\
Max per-block ratio (dimensionless) & 0.820 & \textbf{1.537} \\
Blocks below $1\times$ & 100.0\% & \textbf{98.7\%} (154/156) \\
\bottomrule
\end{tabular}
\end{center}
\textbf{This is a substantive finding, not a cosmetic one.} The defensible claim: $\sim$22--29\% average coverage, $\sim$99--100\% of sampled blocks below $1\times$, at the real observed node count ($\geq$32K).

\textbf{Joint Monte Carlo on the headline} (\texttt{research/sccr\_monte\_carlo.py}, 10{,}000 samples, $N \sim \mathrm{Tri}(10K,150K,\text{mode }60K)$, $C \sim \mathrm{Tri}(\$500,\$2000,\text{mode }\$925)$, $T \sim \mathrm{Tri}(5,30,\text{mode }10)$, $P \sim \mathrm{Tri}(\$30K,\$120K,\text{mode }\$62.9K)$):
\begin{center}
\begin{tabular}{lc}
\toprule
Quantile & SCCR (dimensionless) \\
\midrule
P5 & 0.034 \\
P25 & 0.062 \\
P50 & 0.099 \\
P75 & 0.163 \\
P95 & 0.338 \\
Share below $1\times$ & \textbf{99.8\%} \\
\bottomrule
\end{tabular}
\end{center}
The median is below the deterministic point estimate because the node-count distribution's right tail (median N $\approx$ 86K) dominates. The result is robust in distribution.

\section{Internal Validation, Correction, and Reconciliation (v2.0.0)}\label{sec:6}
Internal validation identified an inconsistency in the SCCR implementation, traced it to a \textbf{duplicated time-horizon term}, corrected the implementation, regenerated all reported values, and confirmed the qualitative conclusions unchanged. Four labeled steps: \textbf{Bug $\rightarrow$ Fix $\rightarrow$ Results $\rightarrow$ Reconciliation}.

\subsection{The bug}\label{sec:6.1}
\texttt{methodology.json} v1.0.0 and \texttt{storage-ratio.js} applied the horizon $T$ twice --- once dividing the denominator of \texttt{costPerBytePerYear} ($C / (B_{\mathrm{all,yr}} / T)$) and again in the lifetime-cost product ($\mathrm{bytes} \times c_b \times T$). This inflated modeled storage cost (USD/block) --- and deflated the ratio --- by exactly \textbf{$10\times$}. The error affected the magnitude, not the logical structure of the model.

\subsection{The fix}\label{sec:6.2}
$c_b$ is defined horizon-free ($C / \text{bytes-per-year}$); $T$ enters only through $L = c_b \times B \times T$. This is the canonical v2.0.0 definition, pinned in \texttt{research/model-spec.json}; no script may reintroduce $T$ into $c_b$.

\subsection{Results}\label{sec:6.3}
\begin{center}
\renewcommand{\arraystretch}{1.2}
\begin{tabular}{lcc}
\toprule
Quantity & Before (v1.0.0 formula, N=60K) & After (v2.0.0, N=60K) \\
\midrule
Cost per byte per year (USD/(byte$\cdot$yr)) & $1.17247\times10^{-7}$ & \textbf{$1.17247\times10^{-8}$} \\
Lifetime storage cost/node/block (USD/block) & \$1.759 & \textbf{\$0.176} \\
Network lifetime cost/block (USD/block) & \$105{,}522 & \textbf{\$10{,}552} \\
\textbf{Average SCCR (dimensionless)} & 0.0172 & \textbf{0.1719} \\
Blocks below $1\times$ & 100\% & \textbf{100\% (unchanged)} \\
\bottomrule
\end{tabular}
\end{center}
\textbf{Why the conclusion survived the correction.} Although the correction increased the estimated SCCR by $\sim10\times$ on a same-capture basis (and $\sim$11.5$\times$ across the as-reported headlines v1.0.0 \textbf{0.0149} $\rightarrow$ v2.0.0 \textbf{0.1719}), every sampled block still remained below the modeled full-cost threshold ($1\times$) under the then-assumed node count (N=60K). At the real N=32K census the average rose further to $\sim$0.225 on the 2026-08-02 live re-measure (168 blocks), with $\sim$99\% of sampled blocks still below $1\times$. The magnitude moved by an order of magnitude; the \emph{direction} of the finding did not.

\subsection{Reconciliation}\label{sec:6.4}
The two cost models (\texttt{storage-ratio.js} and \texttt{utxo\_cost\_model.py}) disagreed by $16.4\times$ (dimensionless). This decomposed as $164\times$ (dimensionless) denominator gap $\div$ $10\times$ (dimensionless) bug $=$ $16.4\times$. The models measure different quantities (block-average vs.\ marginal-inscription attribution); the gap is documented, not erased. See \texttt{verification\_appendix.md}.

\subsection{Reproducibility}\label{sec:6.5}
Every quantity in this paper is regenerated from \texttt{research/model-spec.json} (v2.0.1) by three independent implementations (JS, Python, standalone C); no script redefines a model constant, and the full capture log is retained in \texttt{captured-data/bsahi.db}. A frozen input capture and a cross-check script ship in \texttt{research/reproduce/} (see README).

\section{Limitations and Future Work}\label{sec:7}
\textbf{Limitations:}
\begin{enumerate}[nosep]
  \item The node count ($\geq$32K from the live census) is a \textbf{lower bound}, not a complete enumeration. SCCR is inversely proportional to N: at the live baseline the average inverts above $1\times$ only below $\sim$7.1K nodes, and the ``100\% below $1\times$'' claim breaks below $\sim$49K nodes at the dated capture (\S\ref{sec:5.4}).
  \item Node costs are homogeneous in the model; hardware/bandwidth/electricity vary by geography and operator.
  \item 10-year horizon is an assumption; pruning shortens actual retention, permanent storage extends it.
  \item \textbf{No discounting.} A one-time fee (USD/block) is compared against an undiscounted 10-yr storage-cost sum; discounting the liability at $r=5\%$/yr ($8\%$/yr) reduces the present value by $\sim$27\% (45\%).
  \item Bandwidth is included in the fixed node cost ($C$) yet marginal bandwidth-propagation cost is excluded from the storage leg (fixed-vs-marginal distinction).
  \item Marginal vs.\ average attribution changes the per-byte cost by $164\times$ (dimensionless) --- the choice is explicit and documented.
\end{enumerate}
\textbf{Future work (the resource-pricing program):} UTXO leg (\texttt{utxo\_size\_inc}); bandwidth leg (analytical bounds); validation leg (Core benchmarks); node distribution (\texttt{node\_geo}); BIP-110 pre/post measurement protocol if ever signaled; v3.0 economic-dynamics program (\S\ref{sec:10}).

\subsection{What would falsify this framework?}\label{sec:7.1}
\textit{(Added 2026-08-03, post-advisor review --- see
\texttt{docs/decisions/2026-08-02-publication-decisions.md}.)} A framework that
cannot specify its own failure conditions is not a framework; it is a posture. We
state, in advance, the observations that would force us to retract or
substantially revise the claims in this paper. They operate at two levels --- the
\textbf{measurement} (SCCR, this paper) and the \textbf{framework} (the RIR
family and the ``no single resource market'' thesis; outline in
\texttt{research/framework-paper-outline.md}).

\textbf{Falsifiers of the SCCR measurement:}
\begin{enumerate}[nosep]
  \item \textbf{Independent implementations cannot reproduce the ratio.} The framework ships a frozen capture plus three independent implementations (JS/Python/C, verified agreeing per-block). If an uninvolved party running the protocol on the same inputs cannot reproduce SCCR $\approx$ 0.22 (or the discrepancy cannot be reconciled), the measurement is not trustworthy and the paper's headline falls. This is the external-reproduction critical path (\texttt{research/reproduce/external-reproduction.md}).
  \item \textbf{A better storage-cost model reverses the conclusion.} The denominator ($C \cdot T \cdot N / R_{\mathrm{blocks}}$) rests on a homogeneous node-cost model, an assumed $T=10$ horizon, and the $\geq$32K census. If a defensible refinement --- measured pruned-vs-archival retention, regional cost heterogeneity, or corrected discounting --- moves the banded result ($\sim$22--29\% coverage, most blocks below $1\times$) to a qualitatively different conclusion, the paper's central claim is falsified in its current form.
  \item \textbf{Fees consistently exceed modeled long-term costs.} If sustained fee regimes (or re-measured 2017--2024 fee-peak years) show fees \emph{covering} modeled costs in a stationary way, the partial-internalization reading collapses into ``the market internalizes storage when it matters'' --- a different result, and the banded headline (\S\ref{sec:5}) would be wrong.
\end{enumerate}

\textbf{Falsifiers of the framework (the RIR family and the ``no single resource market'' thesis):}
\begin{enumerate}[nosep]
  \setcounter{enumi}{3}
  \item \textbf{Resource-cost attribution is shown economically inappropriate.} If the planned attribute-pricing regression (working-paper.md \S 11 Q2) shows the single fee price carries \emph{no} persistence signal, then per-resource internalization ratios are category errors, not measurements --- the RIR family premise fails even though SCCR-as-computed may stand.
  \item \textbf{The pruned-vs-archival census shows the storage burden is avoidable at scale.} The $T=10$ replicated-storage denominator assumes archival retention across the network. A measured pruning distribution showing most nodes avoid most of the burden (companion note \texttt{archival-vs-pruned-note.md}) would reframe SCCR as an upper bound on an avoidable cost --- a genuine falsification of the externality reading, though not of the measurement.
  \item \textbf{Measured response functions close the dynamic loop at or above $1\times$.} If node-entry/exit and fee-demand response functions (working-paper.md \S 11 Q1) are measured and exhibit a stable fixed point at or above $1\times$, the persistent-partial-internalization equilibrium hypothesis is falsified.
\end{enumerate}

We do not believe any of these falsifiers currently obtain; we list them so the
reader can check, and so the framework is never mistaken for an unfalsifiable
claim.

\section{Economics and Related Work}\label{sec:8}
\subsection{Is this an externality?}\label{sec:8.1}
Following standard environmental-economics usage (Pigou, 1920; Coase, 1960), a \textbf{negative externality} arises when a transaction's cost is borne by parties who did not consent. Applied to Bitcoin: an inscription's fee is paid by its creator; the long-term storage cost is borne by node operators who neither created the transaction nor were compensated for it. Two qualifications are stated honestly:
\begin{enumerate}[nosep]
  \item \textbf{Voluntary participation.} Node operators choose to run nodes (and may prune). The Pigouvian case is therefore weaker than for a physical externality --- the correct framing is an \emph{unpriced but avoidable} cost, not an imposed one.
  \item \textbf{Measurement, not pricing.} This paper \emph{measures} a ratio; it contains no price mechanism and proposes no policy.
\end{enumerate}
\subsection{Related work and novelty}\label{sec:8.2}
\begin{itemize}[nosep]
  \item \textbf{Aronoff, Praizner, Sabouri (2026), arXiv:2604.17183} --- structural VCG fee model; treats the mempool as a market for scarce block space; does not model storage externalities.
  \item \textbf{Liu, Fang, Cheung, Cai, Huang (2021), arXiv:2103.05866} --- \emph{``An Incentive Mechanism for Sustainable Blockchain Storage.''} Argues storage costs have ``in general not been properly compensated by the users' transaction fees.'' \textbf{This is the closest prior work} and we acknowledge it directly: our contribution is not the observation (already argued in 2021), but a \emph{measured, reproducible, Bitcoin-live-data quantification} (the SCCR) with a reconciliation of cost models and a knife-edge sensitivity bound. \textbf{We do not claim the observation is novel; we claim the measurement is.}
  \item \textbf{Ethereum state-rent / gas-as-state-pricing} and \textbf{Sompolinsky--Zohar} are adjacent threads we build toward but do not model.
\end{itemize}
\subsection{Terminology}\label{sec:8.3}
``Storage Cost Coverage Ratio'' may read as a solvency/insurance term. The economics-native phrasing is \textbf{Cost Internalization Ratio}. We keep SCCR as the operationalized estimator throughout for continuity, but note it is \emph{not} a coverage ratio in the insurance sense --- it is a fee-to-marginal-social-cost comparison.

\section{Conclusion}\label{sec:9}
Whether the fee market internalizes long-lived resource costs is an \textbf{open empirical question}, and this paper contributes the first reproducible measurement using a primary-source node census: a storage-cost-internalization ratio averaging \textbf{$\sim$0.22--0.29 (dimensionless)} (167--156 blocks, N=32K real census, dated captures), with \textbf{$\sim$99--100\% of sampled blocks below the $1\times$ threshold}. The strong form of the claim (``100\% below $1\times$'') does \textbf{not} survive the real node census on the dated capture. The defensible statement is a \textbf{banded estimate}: fees currently cover roughly \textbf{one-fifth to one-third} of the modeled 10-year replicated storage cost depending on the node-count assumption and capture window, with the \emph{direction} (most blocks below $1\times$) robust to every correction made in this paper's review.

\section{Next Evolution (v3.0): When Would the Fee Market Naturally Internalize These Costs?}\label{sec:10}
v2.0 asked and answered a \emph{measurement} question. v3.0 moves from measurement to \emph{economic dynamics}: \textbf{``Under what economic conditions would the fee market naturally internalize these costs?''} The eight-question agenda (preliminary computations at the live fee level, baseline SCCR $=$ 0.2228 at N=32K, T=10, C=\$925/yr):
\begin{center}
\renewcommand{\arraystretch}{1.2}
\begin{tabular}{cll}
\toprule
\# & Question & Headline model output (live baseline) \\
\midrule
1 & Storage horizon: SCCR at $T = 5, 10, 20, 30, 50$ yr & 0.446 / 0.223 / 0.111 / 0.074 / 0.045 (inverse-linear) \\
2 & Storage $10\times$ cheaper: $C = \$92.5$/yr & SCCR $=$ 2.228 --- the gap flips sign \\
3 & BTC $= \$500{,}000$ & SCCR $=$ 1.768; average crosses $1\times$ at $\sim$\$283K \\
4 & Fees sustained at 100 sat/vB for 5 yr & SCCR $\approx$ 11.2 (22.4 at $T{=}5$); crosses $1\times$ at $\sim$9 sat/vB \\
5 & Lightning moves 90\% of payments off-chain & Two-sided; net effect open \\
6 & Nodes double / triple: $N = 32K \to 64K \to 128K$ & 0.223 $\to$ 0.111 $\to$ 0.056 (inverse-linear) \\
7 & Can SCCR reach 1 without protocol changes? & Historically yes: above $1\times$ in 2017--2024 fee-peak years \\
8 & What is the equilibrium? & \textbf{Open.} The model measures a ratio; it does not close the loop \\
\bottomrule
\end{tabular}
\end{center}
\textbf{Open-question honesty.} Nothing in this section claims the fee market \emph{will} internalize these costs, nor that it \emph{should}. The claim is narrower: the question is now well-posed, the drivers are identified, and each driver's lever on the ratio is measured.

\begin{thebibliography}{9}
\bibitem{pigou1920} Pigou, A.~C. \emph{The Economics of Welfare}. Macmillan, 1920.
\bibitem{coase1960} Coase, R.~H. ``The Problem of Social Cost.'' \emph{Journal of Law and Economics} 3 (1960): 1--44.
\bibitem{aronoff2026} Aronoff, D., Praizner, J., Sabouri, S. ``A Model and Estimation of the Bitcoin Transaction Fee.'' arXiv:2604.17183, 2026.
\bibitem{liu2021} Liu, Y., Fang, Z., Cheung, M.~H., Cai, W., Huang, J. ``An Incentive Mechanism for Sustainable Blockchain Storage.'' arXiv:2103.05866, 2021.
\bibitem{bip141} BIP~141: Segregated Witness. Bitcoin Improvement Proposal, 2015--2017.
\bibitem{sompolinsky2015} Sompolinsky, Y., Zohar, A. ``Secure High-Rate Transaction Processing in Bitcoin.'' \emph{FC 2015}.
\end{thebibliography}

\vfill
\noindent\footnotesize\textit{Bitcoin Sahi Research Council --- The Bitcoin Block Space Problem (Working Paper v2.1.0). Component docs: problem\_statement \textbullet{} bip141\_analysis \textbullet{} pruning\_externality\_analysis \textbullet{} utxo\_cost\_note \textbullet{} verification\_appendix \textbullet{} history-of-bitcoin. Conversion note: see header comment for conversion status.}

\end{document}
